\section{Antecedentes}

El estudio del estado del arte constituye un requisito metodológico fundamental para contextualizar la problemática de la gestión y trazabilidad de los servicios de soporte técnico. A continuación, se presentan las investigaciones verificadas en el ámbito internacional, organizadas de forma cronológica ascendente bajo un análisis estructurado en tres momentos: contextualización temática, síntesis parafraseada bajo norma APA 7.\textsuperscript{a} con precisión metodológica y tecnológica, y valoración del aporte específico al presente proyecto de grado.

\subsection{Antecedentes internacionales}

\subsubsection{Investigación 1: \enquote{Helpdesk como alternativa de soporte y gestión tecnológica para una empresa tabasqueña} (Ortiz González, Gómez Lemus y Zurita Cruz, 2021)}

En el ámbito internacional latinoamericano, la estandarización de procesos de atención técnica mediante marcos de gobernanza de tecnologías de información y plataformas de mesa de ayuda representa una estrategia consolidada para superar la desorganización operativa y la falta de canales estructurados de comunicación.

Ortiz González, Gómez Lemus y Zurita Cruz (2021) desarrollaron una investigación orientada a diseñar e implementar una mesa de ayuda (\textit{Helpdesk}) basada en las mejores prácticas de la biblioteca ITIL (\textit{Information Technology Infrastructure Library}) para la empresa de servicios informáticos Techfusione, en Tabasco, México. La metodología articuló las fases de estrategia, diseño, transición y operación de servicios, evaluando la calidad del software mediante métricas de funcionalidad, usabilidad y eficiencia sobre la plataforma ManageEngine ServiceDesk Plus. El estudio reportó la factibilidad de estructurar catálogos de servicios, estandarizar tiempos de respuesta por prioridad y formalizar los flujos de atención desde la recepción hasta la posventa.

\paragraph{Aporte al proyecto.}
El valor sustantivo de este antecedente radica en la aplicación del marco ITIL para catalogar los servicios de mantenimiento e identificar los flujos formales de atención técnica. La presente investigación toma como referencia los criterios de clasificación de incidentes y requerimientos para diseñar la arquitectura del software en Tarija, diferenciándose al incorporar un asistente conversacional con inteligencia artificial enfocado en la interacción directa y automatizada con el cliente.

\subsubsection{Investigación 2: \enquote{Desarrollo de una plataforma web para la gestión de órdenes de servicio de la empresa Computech} (Mosquera Perea, Mazo Agudelo y Martínez Herrera, 2025)}

La transición desde registros informales en papel o mensajería desestructurada hacia plataformas de software web centralizadas constituye una tendencia indispensable para garantizar la persistencia del historial de mantenimiento y la trazabilidad de los procesos de reparación.

Mosquera Perea, Mazo Agudelo y Martínez Herrera (2025), en la Universidad Católica de Pereira (Colombia), desarrollaron una aplicación web orientada a la gestión y seguimiento de órdenes de servicio para la empresa de soporte Computech. Construida bajo la arquitectura Modelo-Vista-Controlador (MVC), utilizando el framework PHP Laravel, MySQL, Bootstrap 5 y la metodología ágil Scrum, la plataforma implementó un esquema de seguridad basado en roles (administrador y técnico) y restricciones de integridad que prohíben la eliminación física de registros. Según la evaluación del sistema (p. 56 del documento original), la búsqueda de información e historial de servicios se redujo en un 60\%, optimizando la eficiencia operativa del taller.

\paragraph{Aporte al proyecto.}
Este trabajo aporta de manera directa el modelo entidad-relación para la gestión de órdenes de trabajo (vinculación entre cliente, dispositivo, diagnóstico y estado de reparación) y la regla de negocio de no eliminación para garantizar una trazabilidad inalterable. El presente proyecto adopta esta estructura de persistencia relacional, pero la extiende al integrar un módulo de asistencia con inteligencia artificial para la consulta remota por parte del usuario final.

\subsection{Síntesis conclusiva del estado del arte}

La revisión de los antecedentes internacionales evidencia que la estandarización de los servicios de soporte técnico mediante software web y marcos de trabajo como ITIL puede reducir los tiempos de búsqueda de información y formalizar la relación con el cliente. No obstante, se identifica una oportunidad de innovación al combinar estas arquitecturas relacionales de trazabilidad con asistentes virtuales conversacionales impulsados por inteligencia artificial, constituyendo el núcleo del diseño propuesto para la ciudad de Tarija.
